\chapter{PENUTUP}

\section{Kesimpulan}

Model re-identifikasi penyu berbasis \textit{backbone} Vision Transformer pra-latih DINOv3
dengan pendekatan \textit{metric learning} mampu mengenali individu penyu dari pola kepalanya pada
skenario \textit{open-set}. Konfigurasi acuan memperoleh Rank-1 sebesar 78,3\% dan mAP sebesar
84,9\%, sedangkan konfigurasi kombinasi memperoleh Rank-1 sebesar 90,4\% dan mAP sebesar 93,5\%.
Penerapan \textit{k-reciprocal re-ranking} hanya menambah 0,9 hingga 1,3 poin persentase. Sebagian
besar kegagalan berupa kesalahan pengurutan pada peringkat teratas, bukan kegagalan menemukan
identitas yang benar. Kegagalan tersebut umumnya terjadi pada citra yang buram dan berkontras
rendah.

Model mampu bergeneralisasi ke data dengan kondisi pengambilan citra yang berbeda. Tanpa
pelatihan ulang, model memperoleh Rank-1 sebesar 81,1\% dan Rank-5 sebesar 97,2\% pada dataset
SeaTurtleID2022. Performanya turun 9,3 poin persentase dari hasil \textit{in-distribution}.
Protokol inferensi pada kedua pengujian identik, sehingga penurunan tersebut murni bersumber dari
pergantian domain. Dengan demikian, generalisasi model disertai biaya performa yang terukur dan
tidak katastrofik.

Penanganan identitas tidak terdaftar merupakan keterbatasan utama sistem. Nilai Registered
Rank-1 turun menjadi 55,3\% pada skema \textit{threshold} global dan 56,7\% pada skema
\textit{background class}, dengan \textit{false-reject rate} sekitar 42\%. Aturan keputusan kedua
skema berbeda secara fundamental, tetapi keduanya bermuara pada titik operasi yang nyaris sama.
Penyebabnya terletak pada distribusi skor. Seluruh skor \textit{top-1} terkurung dalam pita sempit
0,85 hingga 0,98, dan distribusi \textit{unregistered} tertanam di dalam distribusi
\textit{registered}, sehingga separabilitasnya hanya moderat dengan ROC-AUC 0,77 hingga 0,79.
Meskipun demikian, identitas yang benar tetap berada pada lima kandidat teratas bahkan pada kasus
kegagalan terparah. Keterbatasan sistem karenanya bersumber dari mekanisme keputusan dan kalibrasi
skor, bukan dari kualitas \textit{embedding}.


\section{Saran}

Berdasarkan keterbatasan yang teridentifikasi, terdapat beberapa saran yang dapat dipertimbangkan
untuk penelitian selanjutnya, yaitu sebagai berikut.

\begin{enumerate}

    \item Pendekatan \textit{one-factor-at-a-time} perlu didukung dengan rancangan faktorial atau
          ablasi mundur dari konfigurasi kombinasi. Sub-aditivitas yang terbukti pada Subbab 4.4.7
          menunjukkan bahwa penggabungan seluruh nilai terbaik tidak menjamin tercapainya optimum.

    \item Eksperimen perlu dijalankan menggunakan beberapa \textit{seed} yang berbeda. Seluruh hasil
          pada penelitian ini berasal dari satu \textit{seed} tunggal. Akibatnya, parameter dengan rentang
          di bawah 2 poin persentase tidak dapat disimpulkan tanpa estimasi variansi antar-\textit{seed}.
          Pelaporan dalam bentuk rata-rata dan simpangan baku akan memisahkan efek yang nyata dari
          fluktuasi acak.

    \item Penyetelan ambang penolakan perlu dilakukan pada himpunan validasi yang
          terpisah dari himpunan evaluasi. Jumlah identitas tidak terdaftar juga perlu
          diperbesar.

    \item Kemiripan antara profil kiri dan profil kanan penyu dapat dieksploitasi sebagaimana
          diusulkan pada penelitian terdahulu. Pemanfaatannya berpotensi menambah jumlah data yang tersedia
          untuk pencocokan tanpa memerlukan pengumpulan citra baru.

    \item Tahap deteksi dan segmentasi kepala dapat diintegrasikan sehingga terbentuk
          \textit{pipeline} \textit{end-to-end}. Penelitian ini mengasumsikan citra masukan telah
          ter-\textit{crop} pada region kepala. Penggabungan tahap deteksi akan menghasilkan alur kerja yang
          dapat langsung diterapkan pada foto lapangan mentah.

\end{enumerate}